% vim: set spell tw=80:

% vim: set spell tw=80 :

\ifdefined\DRAFT
\documentclass[
  stu,
  draftfirst
]{apa7}
\else
\documentclass[stu]{apa7}
\fi
% vim: set spell tw=80:

\def\NUM{7}
\def\MON{November}
\def\DAY{1}
\def\YEAR{2026}


\usepackage{times}
\usepackage{orcidlink}
\usepackage[style=apa]{biblatex}
\addbibresource{bib.bib}

\title{Critical Thinking \#3}
\author{Ben Raz}
%\author{Ben Raz\,\orcidlink{0009-0003-9576-2066}}
\affiliation{{Student, New Jersey Institute of Technology}}
%\authornote{\addORCIDlink{Ben Raz}{0009-0003-9576-2066}}

\course{ENG 101: Introduction to Academic Writing}
\professor{Maryann Aita}
\duedate{\MON{} \DAY, \YEAR}

\long\def\question#1{}

\ifdefined\DRAFT
\long\def\question#1{
  \fbox{
    \parbox{0.9\textwidth}{
      \sffamily
      \textbf{Question} {\em (Only visible in draft)}
      \begin{quote}
        \begin{quote}#1
        \end{quote}
      \end{quote}
    }
  }
}
\fi

\def\latin#1{\it #1}


\begin{document}

\maketitle

Since the 1979 Islamic revolution, Iran and the United States have had an
adversarial relationship. Iran aims to export its Islamic revolution, a goal
that threatens both moderate Arab monarchies, and Israel, both US allies. Iran
acts on this strategy by supporting a wide range of insurgencies and quasi-state
actors such as \textit{Hezbollah}, the Houthis in Yemen, \textit{Hamas}, and
others. All of these groups either fight Iran's principal enemies or fight their
proxies. Many of these proxies have been described as terrorist groups by
numerous international organizations.

In recent decades, Iran sparked significant anxiety when it launched a nuclear
research program\footnote{While Iran sometimes claims the program's purpose is
peaceful power generation, this claim is severely at odds with observed
enrichment levels, which are too high for power generation and only useful for
weapons use. The program is widely understood to be a nuclear weapons program by
basically all international observers.} The prospect of a nuclear Iran is
considered a grave threat to US national security for a couple of reasons.

The first, and by far the least persuasive fear, is that Iran is a doomsday cult
seeking to end the world and will use their nuclear weapons the moment they gain
them. To put it mildly, such a view is not very well supported.

Another, more legitimate, threat that could come from a nuclear Iran would be
the prospect of further nuclear proliferation in the Middle East. Iran's rivals
in the Gulf, such as Saudi Arabia and the UAE could potentially seek to develop
nuclear weapons of their own\footnote{It is worth noting that Israel, one of
Iran's regional rivals, already possesses nuclear weapons}. Such a nuclear arms
race in a region not known for its political stability would increase the risk
of nuclear weapons being deployed, either accidentally, or by rogue actors, a
terrifying concept for the future of human civilization.

In 2015, the ``P5+1'' (an international coalition including all 5 permanent
members of the UN security council and Germany), after a campaign of
debilitating economic sanctions against Iran, negotiated the Joint Comprehensive
Plan of Action (JCPOA) with Iran. The key stipulations of the agreement were as
follows:

\begin{enumerate}
	\item Iran agrees not to refine Uranium to weapons-grade
	\item Iran agrees to host full and regular IAEA (international atomic energy
		agency) inspections of all nuclear sites
	\item The IAEA can access all undeclared sites within a given period
	\item Some sanctions on Iran will be lifted
	\item If any member of the P5+1 suspects Iran to be enriching Uranium,
		sanctions can be instantly snapped back.
	\item The deal will be up for renewal in 10 years.
\end{enumerate}

The agreement was, by all measures, a success\parencite{crs}. While the deal was
in place, the IAEA found no evidence of nuclear enrichment activity, neither did
the intelligence community, which was most certainly closely monitoring Iranian
activity.

While the deal may been a success, at home, trouble was brewing. Perceptions of
the agreement were stained by a pervasive, albeit understandable, mistrust of
Iran. Many believed that the US was being tricked into the deal, while Iran was
still secretly undertaking nuclear activities. It is important to note that there
has not been any evidence that Iran was violating its duties during the period
in which the agreement was in force. However, for many, distrust of the Iranian
regime prevailed.

Another problem for the agreement was Israel's opposition to it. The Israeli
national security establishment feared what could happen if Israel was left
isolated in a Middle East increasingly open to a detente with Iran. Israel took
some measures against the agreement that may be described as extreme. For
example, the Israeli government extensively surveilled the negotiations and
leaked false stories to the press intended to derail them. Famously, Benjamin
Netanyahu, the prime minister of Israel at the time, visited Congress at the
invitation of Congressional Republicans to give a speech railing against the
agreement. Such an event (a foreign leader supporting the American political
opposition) was completely unprecedented and permanently damaged Netanyahu's
relationship with Obama and with the Democratic Party in general. The event also
underscored the degree to which Netanyahu believed Obama was sidelining him and
showcased his desperation to have some impact on the deal.

In 2016, Donald Trump was elected president after a campaign of opposing the
JCPOA and allying with strict Iran hawks. Trump believed that a more aggressive
approach to Iran was in order and that Iran was reneging on its commitments.
Despite a repeated lack of interest from various P5+1 leaders to renegotiate,
Trump decided to eventually withdraw from the agreement.

In 2018, key sanctions waivers were due for renewal. The deadline came and went
without a renewal. Sanctions on Iran were reinstated; The US had effectively
widthrawn from the agreement. While other P5+1 countries and Iran still {\em de
jure} kept to the agreement, it was clear the agreement was {\em de facto} dead.
Soon after, the IAEA reported that Iran had started work on enriching Uranium
for the first time since the JCPOA was signed.

Since then, the US and Israel have carried out direct strikes on Iranian nuclear
facilities in both 2025 and 2026, marking the first ever direct military
conflict between the three countries. However, Iran has built its facilities
deep underground and continued bombing seems like an unsustainable strategy
given the tremendous cost it has incurred on the US army\footnote{How expensive?
According to the Pentagon \parencite{nyt}, the war cost around \$1.8 billion per
day, just in the first week.}. Iran can continue to rebuild nuclear facilities:
it has both the scientific knowledge as well as the political will. In fact,
Iran is closer and more motivated then ever to build a nuclear bomb. 

Knowing all this, what can be taken away from this situation? First it is important
to understand the why and how of Iran's nuclear ambitions. Iran does not pursue
nuclear weapons because Iran wants to destroy the world. Iran pursues nuclear
weapons technology, at great cost, in the interest of deterrence. Nuclear bombs are the
ultimate insurance policy for any country. Once a country obtains them, no other
country will be able to fight it in an existential war because of the threat of
mutually assured destruction. Think about why the United States has not
attempted to overthrow the much weaker military of North Korea. North Korea has
nuclear weapons. However, most countries do not have nuclear weapons. Why?
Because nuclear weapons are extremely expensive and countries who develop them
have to go through an intense period of international isolation and potential
preemptive military action before finally possessing nuclear weapons. The only
regimes who would take that effort are regimes who feel like their existence
depends on it. American politicians have repeatedly claimed regime change as a
goal in Iran. Such proclamations are dangerous becomes the Iranians are
listening. Finally, it is important to note that Iran can easily make Nuclear
weapons if it really wants to. It would be difficult and costly to remove the
ability. Consider Iran's GDP (around \$200 billion) \parencite{gdp:ir}. If a
country like North Korea with a GDP around ten times smaller \parencite{gdp:ko}
can build a nuclear bomb, then Iran surely can.

What is the solution? There is an old adage that to a hammer, every problem looks
like a nail. The United States, with its intense military power, is tempted
to see every problem as a military problem. But the issue of a potential
nuclear Iran is a {\em diplomatic} issue. And it can only be solved by diplomacy.
Iran has repeatedly demonstrated its interest in negotiation. According to all
available evidence, Iran has abided by the terms of the JCPOA. The idea that
Iran's nuclear ambitions can somehow be militarily crushed has always been a
dangerous delusion. Now, it is becoming more and more unsustainable to keep
maintaining.

\printbibliography

\end{document}
